% Paste this environment into the IEEEtran source.
% The preamble must already \input the sibling iscafig.tex.
\begin{figure*}[t]
\centering
\begin{tikzpicture}[iscafig]
  % ----- before: sp = 12, sb = 0, no link -----
  \begin{scope}[local bounding box=before]
    \node[lbl, anchor=west] at (-0.15,1.15) {\textbf{Before fetch, $sp=12$, $sb=0$}};
    \rfnums
    \foreach \i in {0,...,11} {%
      \rfslot{b}{\i}{occ}{$v_{\i}$}%
    }
    \foreach \i in {12,...,15} {%
      \rfslot{b}{\i}{fre}{}%
    }
    \node[lbl, text=iscaRose, anchor=east] at ([xshift=-1.5mm]b0.west) {$sb$};
    \node[lbl, anchor=north west, align=left] at (-0.15,-2.55)
      {physical $0..11$ resident.\ A growth spill copies $v_0$ through $v_7$.};
  \end{scope}

  % ----- after a growth spill of 8 -----
  \begin{scope}[yshift=-4.35cm, local bounding box=after]
    \node[lbl, anchor=west] at (-0.15,1.15)
      {\textbf{After a growth spill, $sp=5$, $sb=7$}};
    \rfnums
    \foreach \i in {0,...,6} {%
      \rfslot{a}{\i}{fre}{}%
    }
    \rfslot{a}{7}{lnk}{link}
    \rfslot{a}{8}{occ}{$v_{8}$}
    \rfslot{a}{9}{occ}{$v_{9}$}
    \rfslot{a}{10}{occ}{$v_{10}$}
    \rfslot{a}{11}{occ}{$v_{11}$}
    \foreach \i in {12,...,15} {%
      \rfslot{a}{\i}{fre}{}%
    }
    \node[lbl, text=iscaRose, anchor=west] at ([xshift=1.4mm]a7.east) {$sb$};
    \node[lbl, anchor=north west] at (-0.15,-2.70)
      {logical 0 is the link; logical $1..4$ are $v_8..v_{11}$};
  \end{scope}

  % ----- chunk -----
  \begin{scope}[xshift=9.15cm, local bounding box=chunk]
    \node[lbl, anchor=west] at (0,0.55) {\textbf{Spill chunk, 8 values, 48\,B}};
    \rec{0}{+0}{fldhot}{6.35cm}{previous address\,[31:4] \quad count\,[3:0]$=8$}
    \rec{-0.58}{+4}{fldacc}{6.35cm}{tags of $v_0..v_7$, eight nibbles}
    \rec{-1.16}{+8}{fldcore}{6.35cm}{payload $v_0$, oldest}
    \rec{-1.74}{+12}{fldcore}{6.35cm}{payload $v_1$}
    \rec{-2.32}{+16}{fldcore}{6.35cm}{payload $v_2$}
    \rec{-2.90}{+20}{fldcore}{6.35cm}{payload $v_3$}
    \rec{-3.48}{+24}{fldcore}{6.35cm}{payload $v_4$}
    \rec{-4.06}{+28}{fldcore}{6.35cm}{payload $v_5$}
    \rec{-4.64}{+32}{fldcore}{6.35cm}{payload $v_6$}
    \rec{-5.22}{+36}{fldcore}{6.35cm}{payload $v_7$}
    \rec{-5.80}{+40}{fldimm}{6.35cm}{pad to a multiple of 16\,B}
    \node[lbl, anchor=west, align=left] at (0,-6.45)
      {Side tag of word 0 holds count\,[5:4].\\
       A later chunk stores the previous block in [31:4].};
  \end{scope}

  \node[hot, anchor=north west, align=left] (linkenc) at (0,-8.05)
    {\texttt{RF\_LINK}: tag \texttt{0000}, kind \texttt{6} in payload\,[3:0],\\
     address $=$ payload\,[31:4]\,$\ll 4$};
  \node[lbl, anchor=north west, align=left, text width=16.4cm] at (0,-9.15)
    {Dashed cells are free.
     A stack that has popped since the previous spill copies 4 values instead of 8.
     Reload runs only when $sp\le 4$ and logical 0 is a link, so this spill is not undone by the next fetch.};
\end{tikzpicture}
\caption{Register-file spill on a stack that has only been growing.
Fetch spills when $sp\ge 12$, which leaves four free entries, enough for a \texttt{SHORT\_STR}.
The eight oldest values move into one 16-byte-aligned chunk, and the slot that becomes logical 0 is an \texttt{RF\_LINK}.
The spill leaves $sp=5$.
Numbers above the cells are physical indices.
$v_0$ is the oldest resident value.}
\label{fig:regfile}
\end{figure*}
